\documentclass[12pt,a4paper]{article}

\usepackage[top=2cm, bottom=2cm, left=2cm, right=2cm]{geometry}
\usepackage{fontspec}
\setmainfont{Times New Roman}
\usepackage{xcolor}
\usepackage{booktabs}
\usepackage{longtable}
\usepackage{amssymb}
\usepackage{enumitem}

\definecolor{errorred}{RGB}{200,0,0}
\definecolor{warncolor}{RGB}{180,100,0}
\definecolor{okgreen}{RGB}{0,128,0}

\begin{document}

\begin{center}
{\Large\bfseries Academic Review Report}\\[0.5em]
{\large Paper 4: The True Cost of Volatility Targeting}\\[0.5em]
{\normalsize Target Journal: Finance Research Letters}\\
{\normalsize Review Date: 2026-04-04}\\
{\normalsize Review Standard: Strict (FRL referee perspective)}
\end{center}

\section*{Review Summary}

\begin{tabular}{ll}
\textcolor{errorred}{HIGH severity (must fix)} & 4 items \\
\textcolor{warncolor}{MEDIUM severity (should fix)} & 5 items \\
Mild (optional) & 4 items \\
Overall assessment & \textbf{Good --- publishable after revision} \\
Word count & 1,933 / 2,500 (567 words available for additions)
\end{tabular}

\section*{Overall Assessment}

The paper proposes a novel and intuitive decomposition of VT costs into opportunity and transaction components. The 91/9 finding is striking and the VoV conditioning result is practically useful. The ``triple moat'' analysis is interesting but may dilute the focus. The paper is well-written, concise, and fits FRL's scope. However, several issues need addressing before submission.

\section*{\textcolor{errorred}{HIGH Severity Issues (Must Fix)}}

\subsection*{H1: Straw Man Risk in Motivation}

\textbf{Location:} Section 1, paragraph 2

\textbf{Problem:} The paper claims the return shortfall is ``typically attributed to transaction costs'' but does not cite any specific paper making this attribution. This creates a straw man risk --- a reviewer may ask ``who says VT is expensive because of TX costs?''

\textbf{Fix:} Add 1-2 specific citations where practitioners or academics attribute VT costs to turnover/TX. Candidates: Harvey et al.\ (2018) Table 5 reports turnover metrics prominently; Cederburg et al.\ (2020) discuss TX costs as a primary concern.

\subsection*{H2: Missing Key Competitor --- Bongaerts et al.\ (2020)}

\textbf{Location:} Section 1 and throughout

\textbf{Problem:} Bongaerts, Kang \& van Dijk (2020, FAJ) ``Conditional Volatility Targeting'' is the most direct competitor --- they also condition VT on a secondary signal. Not citing this paper is a critical omission that any knowledgeable reviewer will catch.

\textbf{Fix:} Add citation and differentiate: ``While Bongaerts et al.\ (2020) condition on expected returns, we condition on VoV and focus on cost decomposition rather than Sharpe optimization.''

\subsection*{H3: Transaction Cost Assumption Not Justified}

\textbf{Location:} Section 2.1, Eq. 3

\textbf{Problem:} $c = 5$ bps per unit weight change is stated without justification. SPY ETF bid-ask spread is $\sim$1 bp; institutional investors using futures face even lower costs. At 5 bps, the paper may overstate direct costs. Paradoxically, this strengthens the paper's conclusion (opportunity cost share would be even higher), but the assumption must be justified.

\textbf{Fix:} Add a sentence: ``We follow [citation] in assuming 5 bps per leg, which is conservative relative to SPY's typical bid-ask spread of 1--2 bps, thereby providing an upper bound on the direct cost component.''

\subsection*{H4: Cross-OOS Weakness}

\textbf{Location:} Section 3.5

\textbf{Problem:} S2 outperforms BH in only 1/4 periods. This is below the typical 3/5 threshold for robustness. A reviewer may argue the VoV conditioning is sample-specific (worked in COVID but not otherwise).

\textbf{Fix:} This is hard to fix without more data. Acknowledge more prominently: ``The limited cross-OOS success rate (1 of 4) suggests VoV conditioning's value is concentrated in extreme events, consistent with its design as tail-risk insurance rather than systematic alpha.'' Consider adding sensitivity analysis (threshold 0.5 to 1.5) in the 567 remaining words.

\section*{\textcolor{warncolor}{MEDIUM Severity Issues (Should Fix)}}

\subsection*{M1: Triple Moat as Third Contribution}

\textbf{Location:} Section 1, third contribution; Section 3.3

\textbf{Problem:} The ``triple moat'' (diversification + rebalancing + gold crisis alpha) is interesting but may be seen as scope creep for a 2500-word FRL letter. A reviewer might say ``this is a second paper.''

\textbf{Fix:} Downgrade from formal contribution to supporting evidence: ``We further show that a structural rebalancing premium of 54 bps explains why VT cannot outperform simple diversification'' (remove ``triple moat'' terminology, keep the number).

\subsection*{M2: VoV Conditioning --- Which Component Matters?}

\textbf{Location:} Section 2.2

\textbf{Problem:} S2 requires both $z^{VoV} > 1.0$ AND VIX rising. The paper doesn't disentangle which condition drives the result.

\textbf{Fix:} Add a footnote or one sentence: ``Unreported results show that the VoV z-score threshold alone accounts for approximately X\% of the cost reduction, with the VIX-rising filter contributing the remainder.''

\subsection*{M3: Missing Liu et al.\ (2019)}

Liu, F., Tang, X., Zhou, G. (2019) ``Volatility targeting and the smooth sailing paradox'' is directly relevant (discusses when VT destroys value). Should be cited in the literature discussion.

\subsection*{M4: Data Source Credibility}

\textbf{Location:} Section 2.3

Using ``Yahoo Finance (via yfinance)'' as a data source is acceptable for FRL but may raise eyebrows. Consider adding: ``Daily prices are consistent with CRSP to within rounding precision.''

\subsection*{M5: Single-Country Concern}

FRL explicitly discourages ``single-country replications.'' While this paper proposes a new framework (not a replication), the data is US-only. Consider adding one sentence about international applicability or future extension.

\section*{Mild Issues (Optional)}

\begin{enumerate}
\item \textbf{L1}: $\bar{r}$ in Eq.\ 1 should be explicitly defined as ``annualized mean daily return'' or similar.
\item \textbf{L2}: ``Triple moat'' is informal jargon. Replace with ``three structural advantages'' in the formal text.
\item \textbf{L3}: The 12/VIX rule needs a citation (Perchet et al.\ 2016 or similar).
\item \textbf{L4}: VVIX should cite the CBOE methodology paper.
\end{enumerate}

\section*{Strengths (What Works Well)}

\begin{enumerate}
\item \textbf{Novelty}: The 91/9 decomposition is genuinely new and striking.
\item \textbf{Clarity}: The paper is exceptionally well-written for its length.
\item \textbf{Practical value}: The VoV conditioning result is immediately actionable.
\item \textbf{Honest limitations}: Cross-OOS weakness is reported transparently.
\item \textbf{Word count}: At 1,933 words, there is room to address all HIGH issues within the 2,500 limit.
\end{enumerate}

\section*{Revision Priority}

\begin{enumerate}
\item \textbf{H2}: Add Bongaerts et al.\ (2020) citation and differentiation
\item \textbf{H1}: Add specific citations for the ``TX cost attribution'' claim
\item \textbf{H3}: Justify the 5 bps assumption
\item \textbf{H4}: Strengthen cross-OOS discussion + add sensitivity analysis
\item \textbf{M1}: Downgrade triple moat from contribution to supporting evidence
\item \textbf{M2}: Add VoV component disentanglement
\end{enumerate}

\end{document}
